\documentclass[11pt]{article}

\usepackage[margin=0.75in]{geometry}
\usepackage{amsmath, amssymb}
\usepackage{tikz}
\usetikzlibrary{matrix,arrows.meta}
\setlength{\parindent}{0pt}

\begin{document}

\section{Field Extension}

Given two fields $F$ and $K$ with $F \subseteq K$, we say $K$ is an extension of $F$ if $F$ is a subfield of $K$. We denote this by $K/F$. In a field etension, we can treat $K$ as a vector space over $F$. The degree of the field extension is the dimension of $K$ as a vector space over $F$, denoted by $[K:F]$. \\

Let $p(x)\in F[x]$ be an irreducible polynomial, we want to construct a field extension $K$ of $F$ such that $p(x)$ has a root in $K$.\\

The idea is to quotient $F[x]$ by the ideal $(p(x))$. Consider the projection function 

\[\pi: F[x] \to F[x]/(p(x))\]

We claim that $\pi(x)$ is a root of $p(x)$ in $F[x]/(p(x))$. Consider 

\[p(\pi(x)) = p(x+ (p(x))) = \sum_{i=0}^n a_i (x+(p(x)))^i = p(x) + (p(x)) = 0\]

Where the arithmetic is done in the quotient space $F[x]/(p(x))$. Therefore, $F[x]/(p(x))$ is a field extension of $F$ and $x$ is a root of $p(x)$ in $F[x]/(p(x))$. This is a field since $(p(x))$ is a maximal ideal.\\

\textbf{Example 1:} Consider $F=\mathbb{R}$ and $p(x)=x^2+1$. The field extension of $F$ that contains a root of $p(x)$ is given by 

\[\mathbb{R}[x]/(x^2+1)\]

The elements of this field are of the form $a+bx$ where $a, b \in \mathbb{R}$. We can perform basic arithmetic operations in this field. For example, addition gives
\[(a+bx) + (c+dx) = (a+c) + (b+d)x\]
Multiplication is more interesting, since $x^2+1=0$, we have
\begin{align*}
  (a+bx) \cdot (c+dx) &= ac + adx + bcx + bdx^2 \\
  &= ac + adx + bcx + bdx(-1) \\
  &= (ac-bd) + (ad+bc)x
\end{align*}
If we identify $x=i$, then this recovers the complex field $\mathbb{C}$. In this case, we see that we can construct the complex field with such field extension. \\

\textbf{Example 2:} Consider $F=\mathbb{Q}$ and $p(x)=x^2-2$. The field extension of $F$ that contains a root of $p(x)$ is given by 
\[\mathbb{Q}[x]/(x^2-2)\]

The elements of this field are of the form $a+bx$ where $a, b \in \mathbb{Q}$. Multiplication is the field is defined as 
\begin{align*}
  (a+bx) \cdot (c+dx) &= ac + adx + bcx + bdx^2 \\
  &= ac + adx + bcx + bdx(2) \\
  &= (ac+2bd) + (ad+bc)x
\end{align*}
We can identify $x=\sqrt{2}$ in this field, we denote this field by $\mathbb{Q}(\sqrt{2})$. Note that we can as well define $x=-\sqrt{2}$ in this field, in which case we get an isomorphic field $\mathbb{Q}(-\sqrt{2})$. \\

\textbf{Example 3:} Consider $F=\mathbb{Q}$ and $p(x)=x^3-2$. The field extension of $F$ that contains a root of $p(x)$ is given by 
\[\mathbb{Q}[x]/(x^3-2)\]

The elements of this field are of the form $a+bx+cx^2$ where $a, b, c \in \mathbb{Q}$. 

From the above two examples, we can observe a few properties.

\begin{enumerate}
  \item If $p(x)$ has degree $n$, then the elements of $K=F[x]/(p(x))$ are of the form 
  \[a_0 + a_1 x + \cdots + a_{n-1} x^{n-1}\]
  where $a_i \in F$. The basis of $K$ as a vector space over $F$ are $\{1, x, x^2, \cdots, x^{n-1}\}$, and therefore $[K:F]=n$.
  \item The roots of $p(x)$ are algebraically indistinguishable. In particular, defining $x$ to be any root of $p(x)$ gives the same field extension. 
\end{enumerate}

So far we have constructed a field extension for any irreducible polynomial $p(x)$. Therefore, we know that given $F$ and $p(x)$ there exists a field extension of $F$ that contains a root of $p(x)$. However, it is not clear whether our construction gives the minimal field extension. \\

We denote $F(\alpha)$ to be the smallest field that contains both $F$ and $\alpha$. We claim that our construction actually is isomorphic to $F(\alpha)$, this is called the field $F$ generated by $\alpha$. \\

\textbf{Theorem 1:} Let $F$ be a field and $p(x)\in F[x]$ be an irreducible polynomial. Suppose $K$ is an extension field of $F$ containing a root $\alpha$ of $p(x)$. Then 
\[F(\alpha) \cong F[x]/(p(x))\]

\textit{Proof}:\\
Consider the homomorphism $\phi: F[x] \to F(\alpha)$ given by 

\[\phi: f(x) \mapsto f(\alpha)\]

Then $p(x)$ is in the kernel of $\phi$. This induces a homomorphism 

\[\phi: F[x]/(p(x)) \to F(\alpha)\]

Since $p(x)$ is irreducible, the quotient ring $F[x]/(p(x))$ is a field. Therefore, $\phi$ is a field isomorphism. The quotient is isomophic to the image of $\phi$, which is a subfield of $F(\alpha)$, containing both $F$ and $\alpha$. By definition, though, this is exactly $F(\alpha)$. This proves the claim. \\

We now prove that different roots of $F$ has the same algebraic properties. To prove this, we prove a slightly more general result, which will turn useful when we study Galois theory. \\

\textbf{Theorem 2 (Extension Theorem):} Let $F, F'$ be field and $\phi$ be an isomorphism between $F$ and $F'$. The map induces a ring isomorphism $\phi: F[x] \to F'[x]$

\[\phi: p(x) \mapsto p'(x)\]

Where $p'(x)$ is defined by applying $\phi$ to the coefficient of $p(x)$. Let $\alpha, \beta$ be roots of $p(x)$ and $p'(x)$ respectively. Then there is an isomorphism 
\[\sigma: F(\alpha) \to F'(\beta)\]
\[\alpha \mapsto \beta\]
Such that $\sigma$ resticted to $F$ is $\phi$. Pictorially, we have 
\[
\begin{array}{cccc}
\sigma: & F(\alpha) & \xrightarrow{\ \ \sim\ \ } & F'(\beta) \\
& \big\vert &  & \big\vert \\
\phi: & F & \xrightarrow{\ \ \sim\ \ } & F'
\end{array}
\]

\textit{Proof}:\\
The isomorphism induces the isomorphism of fields 
\[F[x] / (p(x)) \xrightarrow{\ \ \sim\ \ } F'[x] / (p'(x))\]
By theorem 1, we see that 
\[F(\alpha) \cong F[x] / (p(x)) \sim F'[x] / (p'(x)) \cong F'(\beta)\]
This proves the claim. \\

To see how this imply the algebraic indistinguishability of roots, we apply the above the theorem with $F'=F$ and $\phi=id_F$. Given distinct roots $\alpha, \beta$, we have $F(\alpha) \cong F(\beta)$. 

\section{Algebraic Extensions}

In the previous section we obtain a field extension $K$ by adjoining a root of the polynomial $p(x)$ to $F$. We can now ask the opposite question: given a field extension $K$ of $F$, can we view $K$ as an extension of $F$ by adjoining roots of some polynomials in $F[x]$?\\

Turns out, the answer is yes, any finite field extension can be viewed as a sequence 
\[F=F_0 \subseteq F_1 \subseteq \cdots \subseteq F_n = K\]
such that each $F_{i+1}$ is an extension of $F_i$ by a root of some polynomial in $F_i[x]$. This is called the tower of fields.\\

We say an element $\alpha\in K$ is algebraic over $F$ if it is the root of some polynomial in $F[x]$. Otherwise, we say that $\alpha$ is transcendental over $F$. The field extension $K/F$ is algebraic if every element of $K$ is algebraic over $F$.\\

We start with the following lemma:\\

\textbf{Lemma (Minimal Polynomial):} Let $\alpha$ be algebraic over $F$. Then there exists a unique monic irreducible polynomial $m_{\alpha}(x)$ with $\alpha$ as a root. We call $m_{\alpha}(x)$ the minimal polynomial of $\alpha$ over $F$. Moreover, a polynomial $f(x)$ has $\alpha$ as a root if and only if $m_{\alpha}(x)$ divides $f(x)$.\\

\textit{Proof}:\\
Let $g(x)$ be the polynomial of least degree with $\alpha$ is a root, and by multiplying by a constant we can assume $g(x)$ is monic. If $g(x)$ is not irreducible, then $g(x)=a(x)b(x)$, so $g(\alpha)=a(\alpha)b(\alpha)=0$, which means that either $a(\alpha)=0$ or $b(\alpha)=0$. This contradicts the minimality of the degree of $g(x)$. Therefore, $g(x)$ must be irreducible and hence is the minimal polynomial of $\alpha$.\\

Now let $f(x)$ be any polynomial with $\alpha$ as a root. We can write 
\[f(x)=g(x)q(x)+r(x)\implies f(\alpha)=r(\alpha)=0\]
By minimiality of $g(x)$, $r(x)$ must be the zero polynomial. Therefore, $f(x)$ is divisible by $g(x)$. \\

We say the degree of $\alpha$ is the degree of $m_{\alpha}(x)$. We then natrually have that 
\[F(\alpha) \cong F[x] / (m_{\alpha}(x))\]
And that $[F(\alpha):F]=\deg(m_{\alpha}(x))$.\\

With this we can now prove a toy version of the theorem.\\

\textbf{Theorem:} The element $\alpha$ is algebraic if and only if the field extension $F(\alpha)/F$ is finite.\\

\textit{Proof}:\\
\textit{($\Rightarrow$):} Suppose $\alpha$ is algebraic over $F$. Let $m_{\alpha}(x)$ be its minimal polynomial, then $[F(\alpha):F]=\deg(m_{\alpha}(x))<\infty$.\\

\textit{($\Leftarrow$):} Suppose $n=[F(\alpha):F]<\infty$. Then the elements 
\[1, \alpha, \alpha^2, \cdots, \alpha^{n}\]
Is linearly dependent. Therefore, there exists scalars $a_0, a_1, \cdots, a_n \in F$ not all zero such that 
\[a_0 + a_1 \alpha + a_2 \alpha^2 + \cdots + a_n \alpha^n = 0\]
This means that $\alpha$ is a root of the polynomial 
\[f(x)=a_0 + a_1 x + a_2 x^2 + \cdots + a_n x^n\]
Therefore, $\alpha$ is algebraic over $F$. \\

In particular, from this we see that if an extension $K/F$ is finite, then it must be algebraic. 

\section{Splitting Fields}

A splitting field of a polynomial $p(x)$ is a field extension $K/F$ such that $p(x)$ has all its roots in $K$. The following theorem shows that such a field extension always exists.\\

\textbf{Theorem:} If $p(x)\in F[x]$ then there exists a splitting field $K/F$ of $p(x)$. \\

\textit{Proof}:\\
We perform induction on the degree $n$ of $p(x)$. If $n=1$ then we may take $K=F$. Suppose the statement is true for all polynomials of degree less than $n$. Let $p(x)$ be a polynomial of degree $n$. If $p(x)$ irreducible component all has degree equals $1$ then taking $K=F$ works. Otherwise let $q(x)$ be an irreducible factor of $p(x)$ with $\deg(q(x))>1$. We know there exists an extension containing a root of $q(x)$. The resulting polynomial has degree $n-1$, so by induction there exists a splitting field $K_1/F(\alpha)$ of $p(x)/(x-\alpha)$. Then $K=K_1(\alpha)$ is a splitting field of $p(x)$. (Take interesction argument).



\end{document}